\documentclass[12pt]{article}
\usepackage{enumitem}
\usepackage{geometry}
\usepackage{amsmath, amssymb}
\geometry{margin=1in}
\begin{document}
\begin{center}
Political Science Assessment 1 \
Total Marks: 40 \
Answer all questions.
\end{center}

\section*{Multiple Choice (20 marks)}
\begin{enumerate}[label=\arabic*.]
\item All of the following are true under the system of checks and balances EXCEPT
\begin{enumerate}[label=(\alph*)]
    \item the president's nominees must be approved by the Senate before taking office
    \item two-thirds of the Senate must approve presidentially negotiated treaties
    \item the president can override a congressional veto and pass laws
    \item Congress can override a presidential veto
\end{enumerate}
\item What is the meaning of "armed attack" in Article 51 UN Charter?
\begin{enumerate}[label=(\alph*)]
    \item Armed attack includes all types of armed force
    \item Armed attack includes all high intensity instances of armed force
    \item Armed attack includes terrorist attacks
    \item An "armed attack" gives the right to invade the aggressor State
\end{enumerate}
\item What does it mean for a State to be sovereign?
\begin{enumerate}[label=(\alph*)]
    \item Sovereignty means being recognised by all other States
    \item Sovereignty means to be able to enter into treaties and join the UN
    \item Sovereignty means freedom to determine one's own affairs without external interference
    \item Sovereignty means possessing the right to defend oneself
\end{enumerate}
\item Why has network news coverage become less diverse in recent years?
\begin{enumerate}[label=(\alph*)]
    \item Decreasing liberal bias in the news media
    \item Increasing visibility for individual candidates on TV
    \item Increasing concentration of ownership in the news media
    \item Decreasing cost of political ads on TV
\end{enumerate}
\item How does Nozick answer the criticism of his historical entitlement theory that if the distribution of goods in society is unjust those at the bottom always...
\begin{enumerate}[label=(\alph*)]
    \item It can be remedied by redistribution of wealth.
    \item If each person's holdings are just, then the total distribution of holdings is just.
    \item Historical factors are secondary to moral imperatives.
    \item He has no answer.
\end{enumerate}
\item Which of the following arguments against the 'fair play' argument in support of a duty to obey the law is the most persuasive?
\begin{enumerate}[label=(\alph*)]
    \item Fairness is a relative term.
    \item The legal system is, in fact, unfair.
    \item It sets a bad example.
    \item The law is irrational and ambiguous.
\end{enumerate}
\item According to Austin the science of jurisprudence is concerned with
\begin{enumerate}[label=(\alph*)]
    \item Morality
    \item Positive law
    \item Divine law
    \item Natural law
\end{enumerate}
\item Both the War Powers Act of 1974 and the Budget and Impoundment Control Act of 1974 represent efforts by Congress to limit the powers of the
\begin{enumerate}[label=(\alph*)]
    \item Joint Chiefs of Staff
    \item House Ways and Means Committee
    \item Central Intelligence Agency
    \item president
\end{enumerate}
\item Ruled unconstitutional in 1983, the legislative veto had allowed
\begin{enumerate}[label=(\alph*)]
    \item the executive branch to veto legislation approved by Congress
    \item federal district courts to overturn legislation
    \item the president to veto state laws
    \item Congress to nullify resolutions approved by the executive branch
\end{enumerate}
\item Why is it important to separate the concept of punishment from its justification?
\begin{enumerate}[label=(\alph*)]
    \item Because its justification depends on the concept employed.
    \item Because any definition of punishment should be value-neutral.
    \item Because the concept of punishment is controversial.
    \item Because the practice of punishment is separate from its justification.
\end{enumerate}
\end{enumerate}

\section*{True / False (20 marks)}
\textit{Answer True or False and justify briefly.}
\begin{enumerate}[label=\arabic*.]
\item True or False: The correct answer to 'Which proposition below is the most consistent with what Rawls claims the POP would opt for in respect of 'social primary goods'?' is 'The POP will choose wealth over a compassionate society.'.
\item True or False: The correct answer to 'Which of the following considers states to be the primary actors in international relations?' is 'Liberalism'.
\item True or False: The correct answer to 'Which one of the following doctrines is not associated with Natural Law thinking' is 'Doctrine of bias'.
\item True or False: The correct answer to 'When the Democratic Party pursues liberal social policies, it is most likely to alienate which of its traditional bases?' is 'Southerners'.
\item True or False: The correct answer to 'Unlike in a closed primary election, in an open primary election' is 'voters may register to vote on the day of the election'.
\item True or False: The correct answer to 'Kelsen's theory of law is called pure theory because Kelsen :' is 'Purely discussed jurisprudence only'.
\item True or False: The correct answer to 'The \_\_\_\_\_\_\_\_ School of jurisprudence believes that the law is an aggregate of social traditions and customs that have developed over the...' is 'Historical'.
\item True or False: The correct answer to 'In which year was the seminal Human Development Report published?' is '1994'.
\item True or False: The correct answer to 'Which of the following would occur if Congress were to pass legislation and declare a recess, and the president took no action on the...' is 'A pocket veto'.
\item True or False: The correct answer to 'The enforcement mechanism of the International Covenant on Civil and Political Rights (and Protocol) consists of' is 'A reporting mechanism and right to individual petition'.
\end{enumerate}

\vfill
\noindent\textit{End of Paper}
\end{document}